\documentclass[12pt]{article}
\usepackage[utf8]{inputenc}
\usepackage[T1]{fontenc}
\usepackage{hyperref}
\title{Saafi-Saafi Language Resources for K-12 Curriculum Development}
\author{Compiled by Claude Sonnet 5.5 (claude-sonnet-5-5) for Andr\'as Kornai}
\date{}
\begin{document}
\maketitle

\section{Introduction}
Saafi-Saafi (ISO 639-3 sav; also Safen, Saafi, Safi-Safi) is a Cangin language of Senegal. It has about 200,000 speakers, concentrated around Thi\`es and Dakar. Five dialects are reported: Boukhou, S\'ebikotane, Sindia, Hasab and Diobass. The language is written in a Latin orthography and, less often, in Arabic script. It has a documented community-led literacy effort.

\section{Primary Educational Resources}
The Association pour le D\'eveloppement de la Langue Saafi (ADLAS), based in Pout, is reported to run a model that teaches reading and mathematics in Saafi and helps children transition to French. Funding is reported to come from the Hewlett Foundation. No published textbooks were seen, and no current enrolment figures were found.
\begin{itemize}
\item \url{https://hewlett.org/?p=28984763}
\end{itemize}

\section{Linguistic and Reference Materials}
\begin{itemize}
\item Botne and Pouille (with Rokhaya Seck), \emph{A Saafi-Saafi (S\'ebikotane Variety) and English // French Dictionary}, 122 pages, with grammar notes. Publisher page: \url{https://www.koeppe.de/titel_a-saafi-saafi-sebikotane-variety-and-english-french-dictionary}
\item Wikipedia on the Safen language, which describes an official Latin orthography with 23 consonants and 5 vowels: \url{https://en.wikipedia.org/wiki/Safen_language}
\item Omniglot page on Safen: \url{https://www.omniglot.com/writing/safen.htm}
\item Omniglot chart: \url{https://www.omniglot.com/charts/safen.pdf}
\item Mbodj (1983) is a doctoral dissertation on the Boukhou variety. It was mentioned in search results only and no link was seen.
\end{itemize}

\section{Religious and Community Resources}
Wikipedia on the Saafi people: \url{https://en.wikipedia.org/wiki/Saafi_people}

\section{Digital and Audio Resources}
Wikitongues page for the language code sav: \url{https://wikitongues.org/languages/sav/}

\section{Curriculum Development Context}
Public schools teach in French, which many Saafi children do not understand. ADLAS's bilingual transition model is the most relevant precedent. The 2016 dictionary and the Latin orthography give a base for materials in the S\'ebikotane variety. The dialect differences would need attention when choosing a standard.
\end{document}